% ============================================================
% Research gap survey, version 1 (arXiv-first).
% Methods for ML + differential equations: what each family has
% solved, and what remains open. Grounded in the papers held
% locally in thesis/literature/{pinn,gnn-operators}/pdfs.
% Build: tectonic research-gap.tex   (or /pdf gap)
% ============================================================
\documentclass[11pt, a4paper]{article}

\usepackage[T1]{fontenc}
\usepackage[utf8]{inputenc}
\usepackage{lmodern}
\usepackage{microtype}
\usepackage[margin=27mm]{geometry}
\usepackage{parskip}

\usepackage{amsmath, amssymb}
\usepackage{booktabs}
\usepackage{array}
\usepackage{enumitem}
\setlist{itemsep=3pt, topsep=5pt}
\usepackage[numbers, sort&compress]{natbib}
\usepackage[hidelinks]{hyperref}
\usepackage[capitalise, noabbrev]{cleveref}
\usepackage{titlesec}
\newcommand{\literaturelink}{\href{https://github.com/blockskode/uio-thesis/tree/main/thesis/literature}{\texttt{thesis/literature/}}}
\titleformat{\section}{\large\bfseries}{\thesection}{0.7em}{}
\titlespacing{\section}{0pt}{16pt}{6pt}
\titleformat{\subsection}{\bfseries}{\thesubsection}{0.6em}{}
\titlespacing{\subsection}{0pt}{12pt}{4pt}

\pagestyle{plain}
\newcolumntype{L}[1]{>{\raggedright\arraybackslash}p{#1}}

\title{\vspace{-14mm}\textbf{Graph Neural Networks that Generalise to Unseen
Domains and Boundary Conditions}\\[3mm]
\large Research gap survey, version 1: arXiv sources}
\author{Yapi Donatien Achou\\
\small Computational Science: Physics, University of Oslo\\
\small Prepared with AI assistance for literature retrieval and drafting}
\date{17 August 2026}

\begin{document}
\maketitle

\section{Thesis objective}

A graph neural network that learns physical laws from simulation data and
generalises to unseen boundary and initial conditions and unseen domains.

This is studied at two scales:

\begin{description}[style=unboxed, leftmargin=0pt, itemsep=2pt,
    font=\normalfont\itshape]
  \item[Continuum mechanics.] Porous media flow: Darcy flow, transport,
  poroelasticity.
  \item[Quantum mechanics.] Kohn-Sham density functional theory, the
  self-consistent field problem.
\end{description}

At each scale the model is used in two ways:

\begin{description}[style=unboxed, leftmargin=0pt, itemsep=2pt,
    font=\normalfont\itshape]
  \item[Surrogate mode.] The model replaces the solver, and is evaluated on
  its accuracy for conditions and geometries it has not seen.
  \item[Warm-start mode.] The model gives the solver a good first guess to
  start from. Solvers work by improving a guess repeatedly until the answer is
  accurate enough, so a better starting guess means fewer rounds of work. The
  solver still decides when it is finished, so correctness is guaranteed by the
  solver rather than by the model, and the model is evaluated only on the
  iterations and wall-clock time it saves.
\end{description}

\textbf{What counts as unseen at the quantum scale.} The external potential
plays the role that boundary conditions and coefficients play at the continuum
scale. It is set by the nuclei, which is chemistry, and by the treatment of
the cell, meaning an isolated molecule against a periodic crystal. So unseen
boundary conditions and geometry becomes unseen chemistry and unseen cell
type.

Unseen chemistry and unseen cell type are both stated as open in the
literature surveyed below: chemistry by the six open questions on
foundation-model potentials~\citep{sixquestions}, and transfer to periodic
systems by the limitations section of the solver-aligned initialisation
work~\citep{sail}.

\textbf{One precision.} One trained model per scale, not one model spanning
both. What transfers across scales is the architecture family, the question
being asked, and the evaluation protocol, not the weights.

\section{What this document is}

This is a survey of the methods that apply machine learning to differential
equations, written to find out what is already solved and what is genuinely
still open. It exists to locate the research gap for a master's thesis asking
whether a graph neural network can learn physical laws from simulation data
and generalise to unseen boundary conditions and unseen domains.

The survey covers 32 papers, all available in \literaturelink{}
and grouped one folder per section. They are grouped by what a method does,
not by which architecture it uses, because the same architecture appears in
several roles and the roles are what the argument turns on.

The recurring finding is that every family succeeds inside the distribution it
was trained on, and that the frontier is the same in all of them: unseen
boundary conditions, unseen geometries and, at the atomistic scale, unseen
chemistries. \Cref{sec:matrix} states that as a table, and the row marked open
is the thesis contribution.

\section{How to read this survey}

Each method family gets three parts: the idea in a few sentences, what the
family has solved with the papers that showed it, and what remains open.
Paper summaries are deliberately short; the papers themselves are in
\literaturelink{} for deep reading.

Every displayed equation and every quoted number carries a pointer into its
source, and the pointer names the paper as well as the location inside it, for
example eq.~(4) of Raissi et al., part~I~\citep{raissi1}. The intent is that
any single statement here can be checked by opening one named paper at one
named place.

\section{Physics-informed neural networks}
\label{sec:pinn}

\subsection{The idea}

A network $u_\theta(t, x)$ represents one solution of one problem, with
$\theta$ the network weights, $t$ time and $x$ the space coordinate. The
differential equation itself becomes the loss. Write the equation as eq.~(1)
of Raissi, Perdikaris and Karniadakis, part~I~\citep{raissi1} does,
%
\begin{equation}
  u_t + \mathcal{N}[u] = 0,
  \qquad x \in \Omega, \quad t \in [0, T],
  \label{eq:pinnpde}
\end{equation}
%
where $u_t$ is the time derivative of the solution, $\mathcal{N}[\cdot]$ is
the nonlinear differential operator that encodes the physics, $\Omega
\subset \mathbb{R}^D$ is the spatial domain in $D$ dimensions and $T$ is the
final time.

The residual $f(t, x) := \partial_t u_\theta + \mathcal{N}[u_\theta]$, which
is eq.~(2) of the same paper~\citep[part~I]{raissi1}, is evaluated by
automatic differentiation, meaning the same machinery that computes gradients
for training is reused to differentiate the network with respect to its input
coordinates. Training then minimises the loss written as eq.~(4) of Raissi
et al., part~I~\citep{raissi1},
%
\begin{equation}
  \mathcal{L}(\theta)
  = \underbrace{\frac{1}{N_u}\sum_{i=1}^{N_u}
      \bigl| u_\theta(t^i_u, x^i_u) - u^i \bigr|^2}_{\text{initial and
      boundary conditions}}
  \; + \;
  \underbrace{\frac{1}{N_f}\sum_{i=1}^{N_f}
      \bigl| f(t^i_f, x^i_f) \bigr|^2}_{\text{the physics}},
  \label{eq:pinnloss}
\end{equation}
%
where $\{t^i_u, x^i_u, u^i\}$ are the $N_u$ known values on the boundary and
at the initial time, and $\{t^i_f, x^i_f\}$ are the $N_f$ collocation points,
which are locations inside the domain where the residual is checked. Nothing
is solved to produce collocation points, so they carry no solution values.

\subsection{Does this family need simulation data?}

For a forward problem, no. The supervision is \cref{eq:pinnpde} itself, plus
the statement of the boundary and initial conditions, so a PINN can be
trained with no solved instance of the problem at all. This is the property
Raissi et al.\ emphasise, and they frame the residual term of
\cref{eq:pinnloss} as a regulariser that makes learning possible in the
small-data regime~\citep[\S{}1, \S{}2.1]{raissi1}.

Data enters in three other situations. Inverse problems use measurements, in
practice often generated by a classical solver. Many later papers add
labelled solution values to stabilise training, which makes them hybrids.
Validation always needs a reference solution, but that is scoring rather
than training.

The contrast with the rest of this survey is worth stating plainly, because
it locates the thesis question. A PINN puts the physics in the loss, needs
no data, and learns a single instance. The operator and simulator families
of \cref{sec:simulators,sec:operators} put the physics in the data and the
architecture, need a solver to generate that data, and learn a whole family
of problems. Transfer to unseen conditions is a question that can only be
asked of the second kind.

\subsection{What this family has solved}

\begin{description}[style=nextline]
  \item[Raissi, Perdikaris and Karniadakis, part I~\citep{raissi1}]
  Introduced the framework in two variants. The continuous-time model treats
  $t$ as an ordinary network input~\citep[\S{}2]{raissi1}, while the
  discrete-time model places a multi-output network on the stages of a
  general Runge-Kutta scheme with $q$ stages~\citep[\S{}3,
  eqs.~(7)--(11)]{raissi1}, which lets implicit high-order time stepping and
  neural networks be used together. On Burgers' equation with viscosity
  $0.01/\pi$, trained on a few hundred boundary and initial points and
  optimised with L-BFGS, a quasi-Newton full-batch method, the reported
  relative $L_2$ error is $6.7 \times 10^{-4}$, about two orders of
  magnitude better than the Gaussian-process method the same three authors
  had published earlier, as reported in \S{}2.1 and fig.~1 of
  part~I~\citep{raissi1}. A Schr\"odinger example uses $50$
  initial points, $50$ periodic-boundary points and $20\,000$ collocation
  points placed by Latin hypercube sampling, with a network of five layers
  and one hundred neurons per layer~\citep[\S{}2.2]{raissi1}. The authors state
  explicitly that the method is not a replacement for finite elements or
  spectral methods, which have fifty years of robustness behind them, and
  they leave uncertainty quantification open~\citep[\S{}4]{raissi1}.
  \item[Part II~\citep{raissi2}]
  The inverse side, where the coefficients $\lambda$ of the operator are
  unknown and are learned jointly with the solution from scattered
  space-time data, by adding them to the trainable
  parameters~\citep[eq.~(5)]{raissi2}. On Burgers' equation, $2\,000$ points
  and a network of nine layers with twenty neurons per layer recover both
  coefficients to well under one percent on clean data, and the reported
  errors stay in the range of a few percent at one to ten percent added
  noise~\citep[Tables~1--2]{raissi2}. On two-dimensional Navier-Stokes the
  identified equation is $u_t + 0.999(u u_x + v u_y) = -p_x +
  0.01047\,(u_{xx} + u_{yy})$ against a true coefficient pair of $1$ and
  $0.01$, degrading only to $0.998$ and $0.01057$ at one percent noise, and
  the entire pressure field is recovered without a single pressure
  measurement, up to the additive constant that Navier-Stokes leaves
  undetermined~\citep[fig.~3]{raissi2}. A Korteweg-de Vries example recovers
  $u_t + 1.000\,u u_x + 0.0025002\,u_{xxx} = 0$ from two time
  snapshots~\citep[fig.~5]{raissi2}.
  This remains the setting where PINNs are genuinely strong, because sparse
  and noisy data enters the loss naturally, and it is also where simulation
  data does appear, since the training measurements are drawn from a
  high-resolution spectral solution.
  \item[Cuomo et al., survey~\citep{cuomo}]
  An eighty-five page review organised around three questions: which network
  is used, how the physical knowledge is represented, and how it is
  integrated~\citep[\S{}1.1]{cuomo}. It also fixes the vocabulary, noting that
  \emph{physics-informed}, \emph{physics-based}, \emph{physics-guided} and
  \emph{theory-guided} are used interchangeably in the literature. Useful as
  an index rather than a read-through.
\end{description}

\subsection{What remains open}

\begin{description}[style=nextline]
  \item[Krishnapriyan et al.~\citep{pinnfailure}]
  The reference negative result, and the single most useful paper in this
  section. On one-dimensional convection with wave speed $\beta$, the plain
  PINN reaches a relative error of $7.5 \times 10^{-1}$ at $\beta = 20$ and
  $9.6 \times 10^{-1}$ at $\beta = 40$, that is nearly one hundred percent
  error after extensive hyperparameter tuning, on a problem with a closed
  form solution~\citep[\S{}3.1, Table~1]{pinnfailure}. The same happens for
  reaction-diffusion as the diffusion coefficient $\nu$
  grows~\citep[eq.~(10), fig.~2]{pinnfailure}. Crucially they show this is
  not a lack of expressive power in the network; the soft PDE constraint of
  \cref{eq:pinnloss} makes the optimisation landscape ill-conditioned, and
  lowering the weight on the residual term buys a smoother landscape at the
  price of a solution that no longer satisfies the
  physics~\citep[\S{}4]{pinnfailure}. Two remedies help: curriculum
  regularisation, which trains at small $\beta$ and reuses those weights to
  initialise the next $\beta$, cuts the $\beta = 20$ error to $9.8 \times
  10^{-3}$~\citep[Table~1]{pinnfailure}; and sequence-to-sequence learning,
  which predicts one time step at a time, cuts a reaction-diffusion case
  from $5.1 \times 10^{-1}$ to $1.2 \times
  10^{-2}$~\citep[Table~2]{pinnfailure}. Both remedies require knowing in
  advance
  which parameter makes the problem hard, and the second quietly reinstates
  the time marching that PINNs were meant to avoid.
  \item[Wang, Yu and Perdikaris~\citep{wangntk}]
  Supplies the mechanism. They derive the neural tangent kernel of a PINN,
  the kernel that describes how an infinitely wide network evolves under
  gradient descent, and prove it converges to a deterministic kernel that
  stays essentially constant during training~\citep[\S{}4,
  Theorem~4.4]{wangntk}. The kernel splits into blocks for the boundary term and
  the residual term, and those blocks have very different eigenvalue scales,
  so the two parts of \cref{eq:pinnloss} converge at very different
  rates~\citep[\S{}7.1, figs.~2--3]{wangntk}. Combined with spectral bias, the
  tendency of such networks to fit smooth components long before sharp ones,
  this explains why high-frequency or multiscale solutions fail. The
  practical output is adaptive loss weighting.
  \item[Wang et al., expert's guide~\citep{expertguide}]
  The state of the art assembled as a pipeline: non-dimensionalise the
  equation, add Fourier feature embeddings, apply random weight
  factorisation, use a modified multilayer perceptron with two input
  encoders, then train with self-adaptive loss balancing, causal training,
  curriculum training and time marching~\citep[fig.~1, \S{}6]{expertguide}. The
  composite loss they minimise separates initial, boundary and residual terms
  explicitly~\citep[eqs.~(2.5)--(2.8)]{expertguide}. The reported accuracies
  document a cliff, shown in \cref{tab:pinnaccuracy}: five decimal digits on
  smooth benchmarks, and roughly ten to twenty percent error as soon as the
  flow becomes multiscale or turbulent, despite multi-GPU training with
  $2 \times 10^5$ iterations per time window on the cylinder
  case~\citep[\S{}7]{expertguide}. For the cylinder case no error is reported at all.
  \item[Zhang et al., limitations review~\citep{pinnlimits}]
  Names three root causes~\citep[abstract]{pinnlimits}: the PDE loss
  approximates multiscale behaviour
  poorly and is ill-conditioned; convergence and error analysis are
  underdeveloped, so the mathematical guarantees that finite elements enjoy
  are missing; and the physics is integrated in a way that leaves the
  residual mismatched with the actual iteration error, meaning a small loss
  does not certify a small error.
\end{description}

\begin{table}[htbp]
  \centering
  \caption{Best reported relative $L_2$ errors from the expert's guide
  pipeline~\citep[Table~1]{expertguide}. The gap between the top and bottom
  groups is
  three orders of magnitude, and it tracks how multiscale the solution is.}
  \label{tab:pinnaccuracy}
  \begin{tabular}{L{62mm} r}
    \toprule
    Benchmark & Relative $L_2$ error \\
    \midrule
    Allen-Cahn equation            & $5.37 \times 10^{-5}$ \\
    Stokes flow                    & $8.04 \times 10^{-5}$ \\
    Advection equation             & $6.88 \times 10^{-4}$ \\
    \midrule
    Lid-driven cavity, $Re = 3200$ & $1.58 \times 10^{-1}$ \\
    Kuramoto-Sivashinsky equation  & $1.61 \times 10^{-1}$ \\
    Navier-Stokes in a torus       & $2.45 \times 10^{-1}$ \\
    Navier-Stokes around a cylinder & not reported \\
    \bottomrule
  \end{tabular}
\end{table}

Two conclusions carry into the rest of this survey. First, accuracy is not
yet competitive with classical solvers on hard problems, and the original
authors say so themselves. Second, and decisive for this project, a PINN is
retrained from scratch for every new instance, so transfer across boundary
conditions or domains is not merely unsolved, it is outside what the
formulation can express.

\section{Learned simulators on graphs}
\label{sec:simulators}

\subsection{The idea}

Represent the discretised system as a graph, with nodes carrying state and
edges connecting neighbours, and learn the time-stepping map itself by
message passing. Physics enters twice, through solver-generated training data
and through the architecture's locality and permutation symmetry. A classical
integrator usually remains, since the network predicts an acceleration or a
rate and a plain Euler step turns that into the next state.

\subsection{What this family has solved}

\begin{description}[style=nextline]
  \item[Battaglia et al.~\citep{battaglia2018}]
  A position paper rather than an experiment, and worth treating as such. It
  defines the graph network block as six functions, three that update edges,
  nodes and a global attribute and three that aggregate them, eq.~(1) of
  Battaglia et al., and shows that interaction networks, message-passing
  networks, graph convolutions and self-attention are all special cases. The
  physics claims it makes are citations to earlier work, not results of its
  own, so this survey uses it for vocabulary and for the argument that
  locality and permutation symmetry are the right inductive biases.
  \item[Sanchez-Gonzalez et al., GNS~\citep{sanchez2020}]
  Particle-based simulation of water, sand and deformable material in one
  framework. Particles are nodes, edges join particles within a fixed radius
  and are rebuilt every step as material moves, the network predicts
  per-particle acceleration, and a semi-implicit Euler integrator advances the
  state. Training is one step at a time, with random-walk noise added to input
  velocities so that rollouts stay stable, a detail without which the method
  does not work (\S{}4.3 of Sanchez-Gonzalez et al.). Generalisation is the
  headline: a model trained on about $2\,500$ particles runs a scene that
  grows to $28\,000$ particles under continuous inflow, and on a domain $32$
  times larger it reaches $85\,000$ particles over $5\,000$ steps, eight times
  the training rollout length (\S{}5.3).
  \item[Pfaff et al., MeshGraphNets~\citep{pfaff2020}]
  The closest published template to this project. The graph carries two edge
  types: mesh edges fixed by the mesh connectivity, and world edges added
  dynamically between nodes that come close in space, which is what lets the
  model handle contact and self-collision, eq.~(1) of Pfaff et al. Removing
  the world edges raises rollout error by 51 and 92 percent on two dynamic
  domains (\S{}5), so the two-edge design is doing real work. Speed is the
  claim that matters commercially: 21 to 23 milliseconds per step against
  $820$ for the ground-truth solver on cylinder flow, and 37 to 38
  milliseconds against $11\,015$ on the aerofoil case, one to two orders of
  magnitude (Table~1 of Pfaff et al.). A model trained on 400-step
  trajectories rolls out stably for $40\,000$ steps (\S{}5).
\end{description}

\subsection{What remains open}

MeshGraphNets is the one paper in this survey with genuine first-party
evidence of mesh topology transfer. A model trained on rectangular cloth is
tested on three disconnected fish-shaped flags and on a windsock with
tassels, whose mesh averages $20\,000$ nodes against about $2\,000$ in
training (\S{}5 of Pfaff et al.). That is a different topology, not a scaled copy
of the training domain, and any honest statement of the geometry gap has to
concede it. The same paper also transfers to unseen physical parameters,
testing aerofoil angles of attack from $-35$ to $35$ degrees after training on
$-25$ to $25$, and Mach numbers from $0.7$ to $0.9$ after training on $0.2$ to
$0.7$, with error rising from $11.5$ to $12.4$ and $13.1$ respectively (\S{}5).

What no paper in this family demonstrates is transfer to a different
\emph{type} of boundary condition. Node types such as walls, obstacles and
inflows are fixed per domain, and none of the three papers holds out a
condition type and tests it unseen.

Cost is the other honest caveat, and the GNS authors report it about their
own method. Their timing appendix shows inference ranging from $51$ percent to $345$ percent of
the ground-truth simulator's step time, meaning it is sometimes slower than
the solver it imitates (appendix~C.6 of Sanchez-Gonzalez et al.). The
authors also state that mesh-based representations are future work, that no
physical symmetries are imposed architecturally, and that efficiency gains
are not yet realised (\S{}6). Pfaff et al.\ report a decoherence effect in cloth
after roughly 50 rollout steps, attributed to chaotic dynamics, and
compensate by reporting a 50-step error alongside the full-trajectory one
(appendix~A.5.2). Both papers depend on a classical solver to generate
thousands of training trajectories, which remains the dominant cost.

\section{Neural operators}
\label{sec:operators}

\subsection{The idea}

Instead of one solution, learn the operator: the map from a problem
specification, meaning coefficients, forcing or initial data, to the solution
function. One trained model then serves a whole family of instances, and the
question of transfer becomes askable, because the model is supposed to handle
inputs it has not seen.

\subsection{What this family has solved}

\begin{description}[style=nextline]
  \item[Lu, Jin and Karniadakis, DeepONet~\citep{deeponet}]
  Gave the field its theoretical footing and a usable architecture. The input
  function is sampled at fixed sensor locations and fed to a branch network,
  the query location goes to a trunk network, and the output is their dot
  product, eq.~(2) of Lu et al., a form justified by the operator
  universal-approximation statement of their eq.~(1). Their Theorem~2 bounds
  the number of sensors needed for a target accuracy on an ODE operator, which
  is rare in this literature. On a diffusion-reaction problem the test error
  reaches about $10^{-5}$ from only 100 input-function samples (\S{}4.3). Note a
  structural consequence of the design: the sensor locations must match
  between training and deployment, so this architecture cannot change
  resolution.
  \item[Li et al., graph kernel networks~\citep{gkn}]
  Made operator learning mesh-independent, and did it with message passing.
  The operator is built from a kernel integral, eq.~(6) of Li et al., whose
  discretisation is exactly a message-passing update over neighbours within a
  radius, their eq.~(10); the kernel is motivated by the Green's function of
  the elliptic problem, their eq.~(5). Because the kernel is a function of
  coordinates rather than of grid indices, one parameter set serves many
  discretisations: trained at $16 \times 16$ and evaluated at $241 \times 241$
  (fig.~1 and Table~1 of Li et al.), with error essentially flat from $85$ to
  $421$ points where a convolutional baseline degrades from $0.0253$ to
  $0.1097$ (Table~8).
  \item[Li et al., Fourier neural operator~\citep{fno}]
  The standard baseline, obtained by restricting the kernel to a convolution
  and parameterising it in Fourier space, eqs.~(3) and~(4) of the FNO paper,
  keeping only the lowest twelve modes for the two-dimensional cases. It is
  the most accurate of the three on shared benchmarks, roughly $0.014$ on
  Burgers across resolutions from $256$ to $8\,192$ points (Table~3) and
  $0.0108$ on Darcy flow against $0.0244$ for the best classical reduced-basis
  competitor (Table~4). Two results give the family its reputation: genuine
  zero-shot super-resolution, training on $64\times64\times20$ data and
  evaluating at $256\times256\times80$ with no retraining (\S{}5.4), and speed,
  $0.005$ seconds per evaluation against $2.2$ for the pseudospectral solver,
  which turns an 18-hour Bayesian inversion into two and a half minutes
  (\S{}5.5). It should be said plainly that the paper contains no theorems; every
  claim is empirical.
  \item[Brandstetter, Worrall and Welling~\citep{brandstetter}]
  A fully neural message-passing solver, with no classical solver at
  inference. Boundary conditions and equation coefficients are supplied to the
  network as an explicit embedding $\theta_{PDE}$ inside the message and
  update functions, eqs.~(8) and~(9) of Brandstetter et al. The ablation on
  that embedding is the most useful number in the paper: on a wave problem
  mixing Dirichlet and Neumann conditions, error is $0.379$ with the
  embedding and $17.591$ without it, roughly forty-six times worse (Table~2).
  Telling the model what its boundary conditions are is therefore not
  optional. Accuracy beats both a fifth-order finite-volume scheme and two FNO
  variants on the harder generalisation setting, $3.70$ against $15.94$ and
  $5.98$ (Table~1), at $0.08$ seconds per rollout against $1.7$ for the
  classical scheme. The paper also contributes the stability techniques,
  temporal bundling and the pushforward trick, that tame autoregressive error.
\end{description}

\subsection{What remains open: the frontier this project targets}

Before the frontier papers, the position of the three foundations should be
stated precisely, because it is easy to overstate. None of DeepONet, the graph
kernel network or FNO varies the domain shape in any experiment: the domains
are an interval, a unit square and a torus. None tests a single trained model
across boundary-condition types, and none tests transfer to a different
governing equation. Only DeepONet reports an out-of-distribution input
experiment at all (fig.~5, \S{}4.1.2 of Lu et al.). FNO argues in a short
paragraph that its architecture is not restricted to periodic conditions, but
supports this by using different fixed conditions in separate experiments
rather than by testing one model across types (\S{}5 of the FNO paper). What this
family solved is resolution and same-distribution inputs, which is a genuine
achievement and a narrower one than the phrase ``learning the operator''
suggests.

\begin{description}[style=nextline]
  \item[Mousavi, Mishra and De Lorenzis~\citep{bcextension}]
  States the boundary-condition gap directly, ICML 2026: the conditioning of
  operators on boundary data ``has received very little attention'', and
  existing methods are restricted to Cartesian geometries, fail to support
  complex conditions, or degrade when conditions vary across samples (\S{}1 of
  Mousavi et al.). Their remedy is an extender module that maps boundary data
  of any type, Dirichlet, Neumann, Robin or mixed, into a latent function
  defined over the whole domain, so that any domain-to-domain operator can
  consume it, eqs.~(4) to~(6). Results are strong: $0.02$ to $0.03$ percent
  error on Poisson against $2.23$ to $3.61$ for the previous
  boundary-specialised baseline (fig.~2a), and elasticity error cut from
  $30.6$ to $11.0$ percent (fig.~2b). Against a finite element reference the
  learned model is comparable in accuracy at $13.8$ against $12.2$ percent and
  roughly twenty thousand times faster, $5.9$ milliseconds against $125$
  seconds (Table~1). Their main backbone, RIGNO, is a message-passing graph
  network.
  \item[Huang et al., Schwarz neural inference~\citep{domaindecomp}]
  States the geometry gap directly, ICLR 2026: existing operators ``lack the
  ability to generalize to entirely novel geometries that differ significantly
  from those present in the training data distribution'', and the need to
  generate new data for each new shape is what blocks industrial use (\S{}1 of
  Huang et al.). Their remedy is local-to-global: train the operator only on
  small simple polygons, then at inference partition an arbitrary domain with
  a classical graph partitioner and stitch subdomain solutions with an
  additive Schwarz iteration, eqs.~(2) and~(3). This is genuine zero-shot
  geometry transfer, including domains with holes, that is a change of
  topology, and a dolphin-shaped domain at $4.5$ percent error. Errors fall
  from 22--28 percent to about $2.1$ percent on Laplace problems and from
  16--167 percent to 5--8 percent on Darcy (Table~1). One detail matters for
  this project: their learned operator is a transformer, not a graph network,
  and their own appendix reports that MeshGraphNets generalises better across
  domains than the transformer they adopted, with replacing the transformer by
  a graph network named as future work (appendix~H.4 and appendix~G).
\end{description}

Now the gap. Three papers sit at this frontier: Mousavi et al.\ and Huang
et al.\ just described, and Brandstetter et al.\ from the previous
subsection. Two things can vary between training and deployment, the type of
boundary condition and the shape of the domain. Each of the three papers
handles one of them and holds the other fixed, and none handles both at once.

Mousavi et al.\ vary boundary conditions richly but fix one geometry per
dataset, and every cross-geometry or cross-condition-type result in that paper
comes from fine-tuning on between 8 and 512 samples of the target setting, not
from zero-shot evaluation (fig.~5 and appendix~F.4). Huang et al.\ achieve
zero-shot transfer to new shapes but never withhold a boundary-condition type;
only condition values vary at test time, and the types are all present in
training (appendix~H.1). Brandstetter et al.\ use a graph network and handle
mixed condition types, but train on a mixture of them rather than holding one
out, and never vary the domain shape at all, so the geometry language in that
paper describes sampling regularity rather than topology.

The untested combination is therefore specific: unseen geometry and unseen
boundary-condition type, varied together in one experiment, on a mesh-based
graph network, evaluated with a solver in the loop so that correctness is
guaranteed rather than hoped for. Nothing in this collection does that, which
is the open row of \cref{tab:matrix}.

Two further observations support attacking it with a graph network rather than
a transformer or a Fourier operator. Structure helps out of distribution:
Shaffer et al.\ report $2.14$ percent against $8.90$ percent on unseen
geometries for an architecture-matched ablation (\cref{sec:structure}). And
Huang et al.\ measured graph networks generalising across domains better than
the transformer they themselves adopted. Neither observation is proof, and both are the kind of
evidence a proposal should cite rather than assert.

\section{Physics built into the architecture}
\label{sec:structure}

\subsection{The idea}

Rather than penalising physics violations in the loss, make them
impossible: choose an architecture whose outputs satisfy the law by
construction. The loss then only has to fit data, because the physics is
already true of every function the network can represent.

\subsection{What this family has solved}

\begin{description}[style=nextline]
  \item[Greydanus, Dzamba and Yosinski, Hamiltonian networks~\citep{hnn}]
  Learn a scalar $H_\theta(q, p)$ and obtain the dynamics as its symplectic
  gradient, following Hamilton's equations, eq.~(2) of Greydanus et al., with
  the network trained on the residual of those equations, their eq.~(3).
  Energy is then conserved by construction rather than approximately. The
  effect is dramatic on the metric that matters: on the mass-spring system
  the energy error falls from $170$ to $0.38$ in the paper's units, and on a
  pendulum learned from pixels from $9.3$ to $0.15$, against an unconstrained
  baseline of the same size (Table~1 of Greydanus et al.).
  \item[Cranmer et al., Lagrangian networks~\citep{lnn}]
  The same trick without requiring canonical coordinates: learn a Lagrangian
  and recover accelerations by solving the Euler-Lagrange equation through
  the network, eq.~(6) of Cranmer et al. This widens the reachable systems,
  and the paper's showcase is a relativistic particle in non-canonical
  coordinates where the Hamiltonian version fails. On a double pendulum the
  energy discrepancy is $0.4$ percent of maximum potential energy against
  $8$ percent for the unconstrained baseline (\S{}5 of Cranmer et al.). The cost
  is a Hessian inverse scaling as $O(d^3)$ in the number of coordinates.
  \item[Shaffer et al., structure-preserving neural PDEs~\citep{structpreserving}]
  The paper that carries this section, because it is the only one of the
  three about PDEs on meshes. It learns a reduced finite-element basis
  conditioned on the mesh together with a discrete flux, then solves the
  resulting nonlinear system at inference, eqs.~(3), (7) and~(9) of Shaffer
  et al. Finite element exterior calculus makes conservation structure and
  Dirichlet conditions exact by construction, and Table~3 of that paper
  reports boundary error of exactly zero against $10^{-3}$ for the
  unconstrained encoder. On standard benchmarks it improves on the best prior
  learned method by 29 percent for elasticity and 71 percent for pipe flow
  (Table~2, \S{}4.1).
\end{description}

\subsection{What remains open}

One caveat about how this survey uses the Hamiltonian and Lagrangian papers.
They are small
conservative systems of ordinary differential equations, a pendulum, a
double pendulum, a two-body problem, and the authors of the Hamiltonian
paper state the restriction themselves: the method assumes a conserved
quantity exists and cannot account for friction (\S{}3.2 of Greydanus et al.).
There are no spatial domains, meshes or boundary conditions in that setting,
so those papers say nothing directly about the transfer axes this survey
cares about. They establish the principle; they do not establish it for
PDEs.

A second caveat concerns the strength of the claim. All three papers compare
a structurally constrained architecture against an unconstrained one, not
against a soft-constrained one, so the defensible statement is that built-in
structure beats no structure. Whether it beats a penalty in the loss is
untested here, and stating otherwise would overreach. The comparison is
cleanest in Shaffer et al., whose baseline shares the same encoder and
differs only by the removal of the structural constraint, which is a real
ablation rather than a different model.

What that ablation shows is the reason this family appears in the survey at
all. On unseen geometries the structured model reaches $2.14$ percent error
against $8.90$ percent for the architecture-matched baseline, and on the
harder decomposition benchmark the coordinate-based baselines fail outright
while the structured model still functions (Table~1 of Shaffer et al.). The
gain is largest exactly where this project is aimed, at out-of-distribution
geometry.

The remaining gaps are stated by Shaffer et al.\ in their limitations
section: two dimensions only, steady state only, a computational mesh and
specified boundary conditions assumed available, extra solve cost against
direct regression, and a learned constitutive model class restricted by
stability requirements. Their boundary-condition transfer evidence is also
weak, being a single qualitative experiment that swaps inlet and outlet
labels without retraining and reports no error figure, while the main
out-of-distribution benchmarks explicitly hold boundary conditions fixed as
geometry varies. Extending exact structure preservation to time-dependent,
dissipative and coupled multiphysics problems is open, and this is where the
project's structured-versus-unstructured ablation sits.

\section{Hybrid methods: accelerate the classical solver}
\label{sec:hybrid}

\subsection{The idea}

Keep the classical solver as the authority and use learning only to reduce
its iteration count, through warm starts, learned preconditioners or learned
solver parameters. Correctness is inherited from the solver, because the
stopping criterion is unchanged, so only speed is at stake. This is the mode
in which a wrong prediction costs time rather than trust.

\subsection{What this family has solved}

\begin{description}[style=nextline]
  \item[Eshaghi et al., NOWS~\citep{nows}]
  The cleanest statement of the warm-start idea. A neural operator trained
  offline maps the problem specification to an approximate solution, and that
  prediction is used only as the initial guess $u_0 = \mathcal{G}_\theta(a)$
  of a Krylov solver, eq.~(5) of Eshaghi et al., which then runs to its usual
  tolerance under the update rule of their eq.~(6). Discretisation,
  preconditioner and stopping criterion are untouched. Reported iteration
  reductions of two- to tenfold and runtime savings of 25 to 90 percent
  (abstract and \S{}1 of Eshaghi et al.); about 90 percent fewer iterations on
  an isogeometric plate with voids across 50 unseen void configurations
  (\S{}3.4, fig.~4c); a Burgers case falling from $9.62$ to $4.75$ seconds on
  average (\S{}3.5, fig.~5b); and a smoke-plume Navier-Stokes case falling from
  $69\,703$ to $17\,438$ total iterations, and from $294.74$ to $93.85$
  seconds, over 150 runs (\S{}3.6). Network inference time is included in the
  reported totals, though the offline data generation and training cost is
  not.
  \item[Zhang, Li and Saad, graph neural multilevel
  preconditioners~\citep{gnnprecond}]
  The benchmarking standard this project must meet, and a genuine graph
  network: message passing within each level of a classical algebraic
  multigrid hierarchy, with interpolation and restriction replaced by learned
  bipartite cross-attention, so that restriction need not be the transpose of
  interpolation, eqs.~(10) and~(17) of Zhang et al. Evaluated on 867
  non-symmetric SuiteSparse matrices from one thousand to one hundred
  thousand unknowns, against algebraic multigrid, incomplete factorisation,
  Jacobi and a single-level graph preconditioner, all under one identical
  solver protocol (\S{}4.1.1). The learned preconditioner wins on iteration
  count in 62.1 percent of cases against the single-level version, rising as
  the tolerance tightens (Table~1), and the multilevel form reaches 72
  percent (Table~2). Failure rates are far below incomplete factorisation,
  which fails to construct on 40.14 percent of these matrices (Table~3).
  Setup cost is disclosed rather than hidden: $210.90$ seconds on average
  against $103.22$ for algebraic multigrid, with the solve phase five times
  faster (Table~4).
  \item[Arisaka and Li, non-intrusive meta-solving~\citep{metasolving}]
  Learning wrapped around a legacy solver that cannot be differentiated at
  all, which is the industrial case. A network predicts solver inputs, and
  since the solver is a black box, its gradient is estimated by a forward
  probe whose variance is cut by a control variate built from a differentiable
  surrogate, eq.~(4) with the variance result of their eq.~(7). Demonstrated
  through PETSc on Poisson, biharmonic and three-dimensional elasticity
  problems: average iterations fall from $73.74$ to $31.54$ on the biharmonic
  case and from $88.71$ to $25.61$ on elasticity (Table~3 of Arisaka and Li),
  and by 74 percent against a supervised-learning baseline for Jacobi
  (Table~2 and \S{}5.1).
  \item[Zou, Xu and Zhang, survey~\citep{itersurvey}]
  Organises the design space into method selection, component optimisation
  and parameter tuning. Two of its conclusions are useful here: the field has
  no standard benchmark datasets, which makes published comparisons hard to
  trust (\S{}5.2 and \S{}5.4 of Zou et al.), and graph networks are named as an
  underexplored direction precisely because a sparse matrix already is a
  graph (\S{}5.4).
\end{description}

\subsection{What remains open}

\begin{description}[style=nextline]
  \item[Wu et al., are hybrid solvers reliable?~\citep{hybridreliable}]
  The cautionary paper this project has to answer. It studies hybrids that
  alternate a classical smoother with a neural correction, and identifies a
  central failure mode named false fixed points: the iteration update shrinks
  to around $10^{-4}$ while the true residual sits at about $10^{1}$, so the
  solver looks converged and is not (fig.~1 and \S{}5.2 of Wu et al.). It also
  prices the fashionable remedy, since training through unrolled iterations
  costs $10.90$ times the training time and $1.61$ times the memory for a
  convergence gain of only $1.16$ to $1.43$ times in one case, and $5.96$
  times the time with $8.53$ times the memory in another (Table~1, \S{}5.4).
  Their own fix, using the physical residual rather than the fixed-point
  update inside Anderson acceleration, converges to $10^{-10}$ on a Helmholtz
  problem where the standard version stalls near $10^{-2}$ (\S{}5.5, fig.~4a),
  with wall-clock gains up to $4.50$ times (Table~2). The honest conclusion
  they draw is that no training objective is universally right and that
  unrolled training should be used with caution (\S{}6).
\end{description}

Reading the five papers of this section together produces the sharpest gap
statement in this survey, and it is a gap of omission rather than of failure.

Not one of the five tests transfer to a different \emph{type} of boundary
condition. NOWS and the reliability study both keep homogeneous Dirichlet
conditions throughout and vary only source terms, coefficient fields or
resolution. The preconditioner paper operates on bare sparse matrices, where
boundary conditions have already been compiled away and no longer exist as a
concept. The meta-solving paper holds the unit domain fixed in every
experiment.

Geometry fares only slightly better. NOWS varies void configurations in one
example (\S{}3.4), and the preconditioner paper generalises across unseen grid
sizes for one equation family, where a single trained model beats
per-matrix-trained baselines at every size (Table~8, appendix~A.4). Nothing
here approaches transfer to a genuinely new domain shape.

Two of the papers name this project's exact next step as their own future
work. Eshaghi et al.\ propose extending warm starts to graph-based operator
architectures for unstructured meshes (\S{}4), and Zou et al.\ argue graph
networks are the natural fit for solver components (\S{}5.4). The combination of
warm starts, graph networks on unstructured meshes, and transfer across
boundary conditions and geometries is asked for in this literature and not
answered by it.

\section{Reinforcement learning for solver control}
\label{sec:rl}

\subsection{The idea}

Reinforcement learning decides where to spend compute rather than what the
solution is. Adaptive mesh refinement is the mature case: an agent looks at
the current mesh, chooses which elements to subdivide, and is rewarded for
error removed per element added. The physics still comes from a conventional
finite element solver, which is called after every refinement step.

\subsection{What this family has solved}

\begin{description}[style=nextline]
  \item[Yang et al.~\citep{rlamr1}]
  Formulates refinement as a Markov decision process over the whole mesh,
  with reward equal to the normalised error reduction from one refinement
  step, eq.~(1) of Yang et al., and the discounted return of eq.~(2). Three
  interchangeable policy heads are compared, a per-element network, a
  hypernetwork and a graph network built from interaction networks, and the
  graph network is the strongest on advection problems, that is on moving
  features (\S{}4.3 and \S{}6.1 of Yang et al.). Reported to match or beat the
  Zienkiewicz-Zhu error estimator on 20 of 24 cases (\S{}6 of Yang et al.), and
  a policy trained on an $8 \times 8$ mesh deploys on $200 \times 200$,
  which is $625$ times more elements (\S{}6.3, fig.~7 of Yang et al.).
  \item[Foucart, Charous and Lermusiaux~\citep{rlamr2}]
  Recasts the same task locally: the agent sees one element at a time and
  chooses to coarsen, do nothing or refine, with a reward that trades
  accuracy against resource cost through a tunable coefficient, eq.~(4) of
  Foucart et al. Because the observation is purely local, the policy is
  independent of mesh size by construction. Reported to preserve solution
  features at under a quarter of the cost of a classical heuristic in
  unsteady one-dimensional advection (\S{}4.3 of Foucart et al.), and to need
  roughly a tenth of the elements on a two-dimensional advecting ring (\S{}4.7
  of Foucart et al.). Notably, \S{}4.2 of the same paper transfers one trained
  policy to a different domain, different boundary conditions and a
  different forcing function, which is a genuine unseen-condition test. The
  policy is a small multilayer perceptron, not a graph network.
  \item[Freymuth et al., swarm refinement~\citep{rlswarm}]
  The most developed of the three, and the only one whose policy is a
  message-passing graph network over the mesh. Every element is an agent
  with a binary refine decision, agents split into child agents, and a
  responsibility mapping propagates credit as elements subdivide, eqs.~(2),
  (4) and (6) of Freymuth et al. Evaluated on seven problem families
  including Stokes flow, linear elasticity and three-dimensional Poisson,
  against uniform refinement, the Zienkiewicz-Zhu estimator, two oracle
  heuristics with privileged access to a fine reference, and
  reimplementations of both papers above. A policy trained on unit-square
  domains deploys on a $20 \times 20$ spiral domain, meshing it in $12.2$
  seconds against about twenty minutes for the uniform reference, roughly a
  hundredfold speedup at a mesh error below $10^{-3}$ (\S{}4.5.1, figs.~15
  and~16 of Freymuth et al.). As few as ten training problems suffice, with
  gains up to a hundred (\S{}4.3.4 of the same paper).
\end{description}

\subsection{What remains open}

Freymuth et al.\ state the binding cost in \S{}5 of their paper: the finite
element system must be solved after every refinement step, both in training
and at deployment, which becomes expensive on fine meshes, and the reward
needs an expensive reference solution during training, so the achievable
resolution is capped by that reference. Their own proposed remedy, learning
to predict refinement regions directly from the problem conditions without
repeated solves, is left as future work. The method is also refinement-only,
stationary, and demonstrated on linear problems.

The other two papers leave related gaps. Yang et al.\ refine one element per
step, exclude coarsening because it needs a two-objective treatment of error
against cost, and describe their explanation for beating greedy baselines as
a conjecture rather than a result (\S{}7 of Yang et al.). Foucart et al.\ flag
that changing the accuracy-versus-cost coefficient requires retraining, with
their tunable-policy workaround described as ongoing research.

Two conclusions matter here. First, transfer is taken seriously in this
family, and Freymuth et al.\ do test unseen domain sizes and unseen material
parameters, although the geometry result relies on a data-augmentation trick
that mimics being a patch of a larger mesh rather than on bare geometric
generalisation. Second, and decisively, in none of the three papers does the
network ever predict a solution field. The graph network chooses where to
refine and a classical solver supplies the physics, so this family is a
control layer over solvers and not a route to learning the physics itself.

\section{The same gap at the atomistic scale}
\label{sec:atomistic}

\subsection{Why this section exists}

Machine-learned interatomic potentials replace density functional theory
inside molecular dynamics, and that community has arrived independently at
the same frontier as the PDE community. The parallel is close enough to be
useful: chemistry is to a potential what boundary conditions and geometry are
to an operator, namely the thing the model was never shown and is expected to
handle anyway.

\subsection{The frontier, stated by the field itself}

\begin{description}[style=nextline]
  \item[Creed et al., six open questions~\citep{sixquestions}]
  A 23-author discussion paper, and the strongest single source for the
  parallel just drawn. It organises the state of universal potentials around six questions
  printed as its section headings: what the minimal definition of an atomistic
  foundation model is (\S{}1), whether the field needs more data, better data or
  better models (\S{}2), whether long-range interactions are handled and whether
  that matters (\S{}3), whether such models can discover genuinely new physics
  (\S{}4), whether they scale to more useful simulations (\S{}5), and how anyone
  knows whether a potential is any good (\S{}6). The verdict in \S{}4 is that models
  learn well within known chemistry and physics but struggle outside their
  training distribution, and the amplification is brutal: energy errors of
  $0.008$~eV can produce migration-barrier errors of $0.74$~eV, two orders of
  magnitude larger (\S{}4.1 of Creed et al.). The paper offers no resolution by
  design.
  \item[Wong and Yang, bias and fine-tuning~\citep{mlipbias}]
  An unseen-chemistry study by construction: a universal MACE potential is
  applied to a liquid electrolyte whose components are absent from its
  training set. The pretrained model reaches $21.24$~meV per atom on an
  independent test set, and iterative fine-tuning brings this to $5.79$, while
  the naive strategy of sampling fine-tuning data with the biased model itself
  plateaus near $10$ and produces unphysical dynamics including spurious
  deprotonation (Tables~1 and~2, fig.~5 of Wong and Yang). The lesson
  transfers directly to this project: data generated by a biased model cannot
  cure that model's bias. Only one architecture was tested, which the authors
  state as a limitation.
  \item[Chiang et al., MLIP Arena~\citep{mliparena}]
  A benchmark platform, NeurIPS 2025 datasets and benchmarks track, that
  evaluates 13 pretrained potentials on physics-grounded tests rather than
  energy regression alone: smoothness of the potential energy surface,
  stability under extreme conditions, and robustness under distribution shift
  measured by a differential-entropy notion of how out-of-distribution a
  structure is (\S{}2.3 of Chiang et al.). Architecturally equivariant models
  keep rotational equivariance exactly, while non-equivariant ones degrade as
  distributional surprise grows, which is direct evidence that built-in
  structure buys robustness. Their related-work section states the field-wide
  position plainly: many models saturate existing metrics yet fail to
  extrapolate to unseen chemistry, strained configurations or
  finite-temperature behaviour.
\end{description}

\subsection{The solver side: learned initial guesses}

\begin{description}[style=nextline]
  \item[Eberhard et al., SAIL~\citep{sail}]
  The atomistic twin of the warm-start idea of \cref{sec:hybrid}, and an
  equivariant graph network. Instead of fitting a converged ground state, it
  learns the starting point of the self-consistent field cycle by
  differentiating through the solver itself, so the objective is alignment
  with the solver's trajectory rather than similarity to a target,
  eqs.~(2)~and~(5) of Eberhard et al. The cycle still runs to convergence, so
  correctness is inherited exactly as in the hybrid family. Reported
  iteration-count reductions of 37, 33 and 28 percent for three functionals on
  molecules up to four times larger than the training set, and a $1.35$-fold
  wall-clock speedup on molecules ten times larger (Table~1 and figs.~1
  and~5). One measurement is worth remembering for cost modelling: building
  the Fock matrix accounts for $97.8$ percent of the work in an iteration
  (fig.~2 of the same paper).
  \item[Hazra, Patil and Sanvito~\citep{densitymatrix}]
  Included as the field's earlier starting point rather than as evidence of
  transfer. A small dense network, not a graph network and not equivariant,
  predicts the one-particle density matrix from atomic coordinates, cutting
  the cycle to one or two iterations and giving a three- to fivefold speedup
  (Table~I and \S{}III.A of Hazra et al.). Equivariance is imposed by hand,
  by fixing molecular orientation, which the authors describe as not general.
  Decisively for this survey, one network is trained per molecule on
  configurations of that same molecule, and no cross-molecule test is
  reported, so the paper makes no transfer claim at all.
\end{description}

\subsection{What remains open}

Eberhard et al.\ list their own boundaries in their limitations section: a
separate model is needed per exchange-correlation functional and, for matrix
representations, per basis set; the evaluation covers only hydrogen, carbon,
nitrogen, oxygen and fluorine; and transfer to heavier elements, to
non-equilibrium geometries and to periodic systems is untested, with the last
described as remaining open. So even the best transferable-initialisation
result is transfer in system size, not in chemistry.

That distinction matters for how this survey uses the analogy. Unseen size is
demonstrated, unseen chemistry is not, which is precisely the pattern found in
\cref{sec:operators}, where resolution transfer is solved and boundary
conditions and geometry are not. The two scales are not merely similar in
spirit, they are stuck on the same step.

\section{Synthesis: the generalisation matrix}
\label{sec:matrix}

\begin{table}[htbp]
  \centering
  \small
  \begin{tabular}{@{}L{4.6cm}lL{5.6cm}@{}}
    \toprule
    Generalisation axis & Status & Evidence \\
    \midrule
    New forcing or coefficients, within the training distribution &
    largely solved & operator learning~\citep{deeponet, fno} \\
    \addlinespace
    New mesh resolution &
    solved for spectral operators & zero-shot
    super-resolution~\citep{fno}; mesh-invariant kernels~\citep{gkn} \\
    \addlinespace
    Larger systems, longer rollouts &
    partly solved & particle and mesh simulators~\citep{sanchez2020,
    pfaff2020}, with stability tricks~\citep{brandstetter} \\
    \addlinespace
    \textbf{Unseen boundary-condition types} &
    \textbf{open} & gap stated in print, and closed only by fine-tuning on
    the target setting~\citep{bcextension}; mixed types trained as a
    mixture, never held out~\citep{brandstetter} \\
    \addlinespace
    \textbf{Unseen geometries} &
    \textbf{partly open} & mesh topology transfer demonstrated for
    cloth~\citep{pfaff2020}; zero-shot for shapes with holes, but with a
    transformer, not a graph network~\citep{domaindecomp}; structure helps
    out of distribution~\citep{structpreserving} \\
    \addlinespace
    \textbf{Both together, in one experiment} &
    \textbf{open} & each frontier paper closes one axis and fixes the
    other~\citep{bcextension, domaindecomp, brandstetter} \\
    \addlinespace
    \textbf{Unseen chemistries (atomistic)} &
    \textbf{open} & field-stated open questions~\citep{sixquestions};
    bias on unseen chemistry~\citep{mlipbias}; benchmark
    evidence~\citep{mliparena} \\
    \addlinespace
    A different governing equation &
    far open & no credible general result to date \\
    \bottomrule
  \end{tabular}
  \caption{What is solved and what is open, per generalisation axis. The
  bold rows are the research gap this project addresses; the last row is
  deliberately out of scope.}
  \label{tab:matrix}
\end{table}

Three observations close the survey.

Every family performs inside its training distribution, and the frontier
is the same everywhere: the axes along which real applications vary, since
every patient, reservoir and molecule is a new boundary condition,
geometry or chemistry.

The most promising lever reported so far is structure in the architecture
rather than volume of data~\citep{structpreserving, hnn}, which is a
testable hypothesis rather than a folklore claim.

And the hybrid deployment mode~\citep{nows, gnnprecond} changes what
failure costs: outside the distribution a surrogate is silently wrong,
while a warm start merely saves fewer iterations. Measuring where those
two behaviours separate is precisely the proposal built on this survey.

\begin{thebibliography}{99}
\small

\bibitem{raissi1} M.~Raissi, P.~Perdikaris and G.~E. Karniadakis. Physics
informed deep learning (part I): data-driven solutions of nonlinear partial
differential equations. arXiv:1711.10561, 2017.

\bibitem{raissi2} M.~Raissi, P.~Perdikaris and G.~E. Karniadakis. Physics
informed deep learning (part II): data-driven discovery of nonlinear
partial differential equations. arXiv:1711.10566, 2017.

\bibitem{cuomo} S.~Cuomo et al. Scientific machine learning through
physics-informed neural networks: where we are and what's next.
arXiv:2201.05624, 2022.

\bibitem{pinnfailure} A.~S. Krishnapriyan et al. Characterizing possible
failure modes in physics-informed neural networks. NeurIPS 2021; arXiv:2109.01050.

\bibitem{wangntk} S.~Wang, X.~Yu and P.~Perdikaris. When and why PINNs
fail to train: a neural tangent kernel perspective. arXiv:2007.14527, 2020.

\bibitem{expertguide} S.~Wang et al. An expert's guide to training
physics-informed neural networks. arXiv:2308.08468, 2023.

\bibitem{pinnlimits} W.~Zhang et al. Physics-informed neural networks as
intelligent computing technique for solving PDEs: limitation and future
prospects. arXiv:2411.18240, 2024.

\bibitem{battaglia2018} P.~W. Battaglia et al. Relational inductive biases,
deep learning, and graph networks. arXiv:1806.01261v3, October 2018. Position paper; no first-party experiments.

\bibitem{sanchez2020} A.~Sanchez-Gonzalez et al. Learning to simulate
complex physics with graph networks. ICML 2020, PMLR vol.~119; arXiv:2002.09405.

\bibitem{pfaff2020} T.~Pfaff et al. Learning mesh-based simulation with
graph networks. ICLR 2021; arXiv:2010.03409v4.

\bibitem{deeponet} L.~Lu, P.~Jin and G.~E. Karniadakis. DeepONet: learning
nonlinear operators based on the universal approximation theorem of
operators. arXiv:1910.03193v3, April 2020.

\bibitem{gkn} Z.~Li et al. Neural operator: graph kernel network for
partial differential equations. arXiv:2003.03485, March 2020.

\bibitem{fno} Z.~Li et al. Fourier neural operator for parametric partial
differential equations. ICLR 2021; arXiv:2010.08895v3.

\bibitem{brandstetter} J.~Brandstetter, D.~E. Worrall and M.~Welling.
Message passing neural PDE solvers. ICLR 2022; arXiv:2202.03376v3.

\bibitem{hnn} S.~Greydanus, M.~Dzamba and J.~Yosinski. Hamiltonian neural
networks. NeurIPS 2019; arXiv:1906.01563v3.

\bibitem{lnn} M.~Cranmer et al. Lagrangian neural networks.
arXiv:2003.04630v2, July 2020.

\bibitem{bcextension} S.~Mousavi, S.~Mishra and L.~De Lorenzis. Imposing
boundary conditions on neural operators via learned function extensions.
ICML 2026, PMLR vol.~306; arXiv:2602.04923v2, May 2026.

\bibitem{domaindecomp} J.~Huang et al. Operator learning with domain
decomposition for geometry generalization in PDE solving. ICLR 2026; arXiv:2504.00510v2, February 2026.

\bibitem{structpreserving} B.~D. Shaffer, S.~Koohy, B.~Kinch,
M.~A.~Hsieh and N.~Trask. Structure-preserving learning improves geometry
generalization in neural PDEs. ICML 2026, PMLR vol.~306; arXiv:2602.02788v2,
June 2026.

\bibitem{nows} M.~S.~Eshaghi, C.~Anitescu, N.~Valizadeh et al. NOWS:
neural operator warm starts for accelerating iterative solvers.
arXiv:2511.02481v4, May 2026.

\bibitem{gnnprecond} Z.~Zhang, R.~P.~Li and Y.~Saad. Graph neural
multilevel preconditioners for iterative solvers. KDD~'26;
arXiv:2607.28456, July 2026.

\bibitem{metasolving} S.~Arisaka and Q.~Li. Accelerating legacy numerical
solvers by non-intrusive gradient-based meta-solving. ICML 2024, PMLR
vol.~235; arXiv:2405.02952.

\bibitem{itersurvey} H.~Zou, X.~Xu and C.-S.~Zhang. A survey on
intelligent iterative methods for solving sparse linear algebraic equations.
arXiv:2310.06630, 2023.

\bibitem{hybridreliable} Y.~Wu, J.~W.~van Beek, V.~Dolean and
A.~Heinlein. Are deep learning based hybrid PDE solvers reliable? Why
training paradigms and update strategies matter. arXiv:2602.06842v2, June
2026.

\bibitem{rlamr1} J.~Yang, T.~Dzanic, B.~Petersen et al. Reinforcement
learning for adaptive mesh refinement. AISTATS 2023, PMLR vol.~206;
arXiv:2103.01342.

\bibitem{rlamr2} C.~Foucart, A.~Charous and P.~F.~J.~Lermusiaux. Deep
reinforcement learning for adaptive mesh refinement. arXiv:2209.12351, 2022.

\bibitem{rlswarm} N.~Freymuth, P.~Dahlinger, T.~W\"urth, S.~Reisch,
L.~K\"arger and G.~Neumann. Adaptive swarm mesh refinement using deep
reinforcement learning with local rewards. arXiv:2406.08440v2, January 2026;
extended version of the earlier ASMR conference paper, method named ASMR++.

\bibitem{sixquestions} I.~Creed, T.~Rein, I.~Vitenburgs et al.\ (23
authors). Six open questions in machine-learned interatomic potential
foundation models. arXiv:2606.07327v2, June 2026.

\bibitem{mlipbias} N.~H.~Wong and J.~H.~Yang. Bias in universal
machine-learned interatomic potentials and its effects on fine-tuning.
arXiv:2603.10159v2, June 2026.

\bibitem{mliparena} Y.~Chiang, T.~Kreiman, C.~Zhang et al. MLIP Arena:
advancing fairness and transparency in machine learning interatomic
potentials via an open, accessible benchmark platform. NeurIPS 2025,
datasets and benchmarks track; arXiv:2509.20630v2.

\bibitem{sail} E.~S.~Eberhard, V.~Kotsev, T.~G\"uthle and
S.~G\"unnemann. Transferable SCF-acceleration through solver-aligned
initialization learning. arXiv:2604.21657v2, May 2026.

\bibitem{densitymatrix} S.~Hazra, U.~Patil and S.~Sanvito. Predicting
the one-particle density matrix with
machine learning. arXiv:2401.06533, 2024.

\end{thebibliography}

\end{document}
